\section{What this does not show}\label{sec:discussion}

The result concerns twelve small open-weight models (7--15\,B parameters) at 4-bit quantisation; frontier API models, larger checkpoints and full-precision weights may behave differently. The eight-market menu, the templates and their translations are one instrument among many; another menu, a request for a single-country recommendation, or native-speaker templates could change the picture. Decoding is greedy, so each estimate describes a model's modal reply, not the distribution a user would sample. Origin is the lead developer's country, which is coarse: Mistral-NeMo was co-developed with NVIDIA, and fine-tunes inherit most of their base model. The outputs say nothing about any developer's intent or training data, and nothing about whether the advice is good: an allocation can be free of home tilt and still be far from any sensible benchmark, as these are.

\section{Conclusion}\label{sec:conclusion}

Open-weight LLMs from five origins do not allocate more to their home equity market when the prompt is held fixed: the pre-registered home effect is \betaHome{}\,pp, and a post-hoc check for dispersion puts it near zero. What moves the allocation is what the user says. The investor's stated residence is associated with \gammaInv{}\,pp more for that market, \gammaRatio{} times the absolute size of the pooled home effect, and every model holds cash in dollars. For users and supervisors, the practical tilt to watch in these models is towards the user's own framing and towards equal weighting, not towards the developer's country.
